\documentclass[12pt]{book}

% ==========================================
% PREÁMBULO: Configuración y Paquetes
% ==========================================

\usepackage{import}
\import{config/}{package}

% -- comando nuevos
\import{config/}{new_comandos}
% Ajustamos la lista para que el texto de la definición baje al siguiente renglón
\setlist[description]{style=nextline, leftmargin=15pt}

% -- Configuración de estilos 
\import{config/}{styles}

% --- Configuración de Entornos Matemáticos ---
% (El contador [chapter] hace que se numeren según el capítulo: 1.1, 1.2, etc.)
\import{config/}{reglas_teoria}
% --- Información del Libro ---
 % O puedes poner un año fijo: "2026"
\title{}
\author{}
\date{}
% ==========================================
% INICIO DEL DOCUMENTO
% ==========================================
\begin{document}
\begin{titlepage}
    % Desactivar los márgenes estándar para la portada usando TikZ
    \begin{tikzpicture}[remember picture, overlay]
        
        % 1. Fondo de la portada
        \fill[bgColor] (current page.south west) rectangle (current page.north east);
        
        % 2. Franja lateral o decorativa (Estilo minimalista/geométrico)
        \fill[primaryColor] (current page.north west) rectangle ($(current page.south west) + (1.5cm, 0)$);
        \fill[accentColor] ($(current page.north west) + (1.5cm, 0)$) rectangle ($(current page.south west) + (1.7cm, 0)$);

        % 3. Bloque de Geometría Estocástica / Elemento Alusivo a Probabilidad
        % Aquí dibujamos una abstracción de distribución normal y redes de probabilidad
        \begin{scope}[shift={($(current page.south east) + (-6cm, 6cm)$)}, scale=1.2]
            % Curva de Gauss de fondo en líneas punteadas
            \draw[very thick, secondaryColor!40, smooth, domain=-3:3, samples=50] 
                plot (\x, {3*exp(-\x*\x/2)});
            \fill[secondaryColor!10, opacity=0.5, domain=-1.5:1.5, samples=50] 
                (-1.5,0) -- plot (\x, {3*exp(-\x*\x/2)}) -- (1.5,0) -- cycle;
            
            % Espacio muestral / Nodos de procesos estocásticos (Caminata aleatoria / Grafo)
            \node[circle, fill=primaryColor, inner sep=2pt] (A) at (-2,1) {};
            \node[circle, fill=accentColor, inner sep=3pt] (B) at (-0.5,2.5) {};
            \node[circle, fill=secondaryColor, inner sep=2.5pt] (C) at (1,1.8) {};
            \node[circle, fill=primaryColor, inner sep=2pt] (D) at (2.5,0.5) {};
            
            \draw[thick, primaryColor!70, ->, >=stealth] (A) -- (B) node[midway, above left, scale=0.7] {$p_1$};
            \draw[thick, primaryColor!70, ->, >=stealth] (B) -- (C) node[midway, above, scale=0.7] {$p_2$};
            \draw[thick, primaryColor!70, ->, >=stealth] (C) -- (D) node[midway, right, scale=0.7] {$1-p_1$};
            \draw[thick, dashed, secondaryColor!60] (A) to[out=45,in=135] (C);
        \end{scope}

        % 4. Textos de la Portada
        % Institución / Sello editorial (Opcional)
        \node[anchor=north west, text=primaryColor, font=\large\scshape\bfseries] 
            at ($(current page.north west) + (3cm, -3cm)$) {Especialización en ciencia y analítica de datos};

        % Título Principal
        \node[anchor=west, text=primaryColor, font=\LARGE\bfseries, scale=1.3] (Title)
            at ($(current page.north west) + (3cm, -7cm)$) {Sistemas de aprendizaje automático};
            
        % Subtítulo o enfoque
        \node[anchor=west, text=secondaryColor, font=\Large\itshape, text width=12cm] (Sub)
            at ($(Title.south west) + (0, -0.8cm)$) {Aprendizaje de datos históricos};

        % Línea divisoria elegante
        \draw[accentColor, ultra thick] ($(Sub.south west) + (0, -0.5cm)$) -- ($(Sub.south west) + (6cm, -0.5cm)$);

        % Nombre del Autor
        \node[anchor=west, text=black!80, font=\Large\scshape] 
            at ($(current page.north west) + (3cm, -12cm)$) {Javier Andrés Causil Martínez };

        % Pie de página (Año / Edición / Ciudad)
        \node[anchor=south west, text=primaryColor!80, font=\small\sffamily] 
            at ($(current page.south west) + (3.5cm, 2cm)$) {Primera Edición $\cdot$ 2026};

    \end{tikzpicture}
\end{titlepage}

\clearpage
% --- Páginas Preliminares ---
\frontmatter 
% Las páginas aquí usarán números romanos (i, ii, iii...)

\maketitle % Genera la portada con la info de arriba

\tableofcontents % Genera el índice general
\listoffigures % Descomenta si quieres un índice de imágenes
\listoftables  % Descomenta si quieres un índice de tablas
\chapter{Prefacio}
Este libro busca ser un compendio de práctica y teoría para entender el aprendizaje automático.

% --- Contenido Principal ---
\mainmatter 
% Las páginas aquí usarán números arábigos (1, 2, 3...)
\part{Algoritmos y Modelos en Machine Learning}
\import{part/part_0/}{main}
\part{Gestión del datos}
\import{part/part_1/}{main}
\part{Estadística}
\import{part/part_2/}{main}


% --- Páginas Finales ---
\backmatter 
% Para apéndices, glosario y bibliografía

% Ejemplo de bibliografía manual (Para libros más grandes se recomienda usar BibTeX o Biber)
\bibliographystyle{plain} 

% 2. Llama a tu archivo .bib (sin escribir la extensión .bib)
\bibliography{referencias/libros}
\end{document}